\SourcePage{30}{35}
\section{Spherical photon waves}

Having determined the possible values of photon angular momentum, we must now find the corresponding wave functions.\footnote{This question was first considered by Heitler (W.~Heitler, 1936). The form of the solution presented here is due to V.~B.~Berestetskii (1947).}

First consider the formal problem of finding vector functions that are eigenfunctions of the operators $\widehat{\mathbf j}^{\,2}$ and $\widehat j_z$. At this stage we neither decide which of these functions enter the photon wave functions of interest nor impose transversality.

We seek the functions in the momentum representation. The coordinate operator in this representation is $\widehat{\mathbf r}=i\partial/\partial\mathbf k$ (see III, (15.12)). The orbital-angular-momentum operator is
\[
\widehat{\mathbf l}=\widehat{\mathbf r}\times\mathbf k
=-i\left(\mathbf k\times\frac{\partial}{\partial\mathbf k}\right),
\]
and thus differs from the angular-momentum operator in the coordinate representation only by replacing $\mathbf r$ with $\mathbf k$. The formal solution of the problem is therefore the same in both representations.

Denote the required eigenfunctions by $\mathbf Y_{jm}$ and call them vector spherical harmonics. They must obey
\begin{equation}
\widehat{\mathbf j}^{\,2}\mathbf Y_{jm}=j(j+1)\mathbf Y_{jm},
\qquad
\widehat j_z\mathbf Y_{jm}=m\mathbf Y_{jm}
\tag{7.1}
\end{equation}
(the $z$ axis is a specified direction in space). We shall show that every function of the form $\mathbf aY_{jm}$ has this property, where $\mathbf a$ is any vector constructed from the unit vector $\mathbf n=\mathbf k/\omega$, and $Y_{jm}$ is an ordinary scalar spherical function. We use throughout the definition from III, \S~28:
\begin{equation}
Y_{lm}(\mathbf n)=(-1)^{(m+|m|)/2}i^l
\sqrt{\frac{(2l+1)(l-|m|)!}{4\pi(l+|m|)!}}\,
P_l^{|m|}(\cos\theta)e^{im\varphi}
\tag{7.2}
\end{equation}
\SourcePage{31}{36}
\noindent
($\theta$ and $\varphi$ are the spherical angles specifying the direction $\mathbf n$).\footnote{For later reference, we record the value at $\theta=0$, when $\mathbf n$ lies along the $z$ axis:
\begin{equation}
Y_{lm}(\mathbf n_z)=i^l\sqrt{\frac{2l+1}{4\pi}}\,\delta_{m0}.
\tag{7.2a}
\end{equation}}
Recall the commutation rule III, (29.4):
\[
\{\widehat l_i,a_k\}_{-}=ie_{ikl}a_l.
\]
The right-hand side can be written $(-\widehat s_i a_k)$, where $\widehat{\mathbf s}$ is the spin-1 operator; its action on a vector function is precisely defined by $\widehat s_i a_k=-ie_{ikl}a_l$ (see III, \S~57, problem~2). Hence
\[
\widehat l_i a_k-a_k\widehat l_i=-\widehat s_i a_k.
\]
Using this equality, we find
\[
\widehat j_i a_k=(\widehat l_i+\widehat s_i)a_k=a_k\widehat l_i.
\]
Consequently,
\[
\widehat{\mathbf j}^{\,2}(\mathbf aY_{jm})
=\mathbf a\widehat{\mathbf l}^{\,2}Y_{jm},
\qquad
\widehat j_z(\mathbf aY_{jm})=\mathbf a\widehat l_zY_{jm}.
\]
But $Y_{jm}$ is an eigenfunction of $\widehat{\mathbf l}^{\,2}$ and $\widehat l_z$ with eigenvalues $j(j+1)$ and $m$, respectively, and equations~(7.1) follow.

Three essentially different types of vector spherical harmonic result from choosing $\mathbf a$ to be one of the following vectors:\footnote{The operator $\boldsymbol\nabla_{\mathbf n}=|\mathbf k|\boldsymbol\nabla_{\mathbf k}$ acts on functions that depend only on the direction $\mathbf n$. In spherical coordinates it has only two components:
\[
\boldsymbol\nabla_{\mathbf n}=\left(\frac{\partial}{\partial\theta},
\frac{1}{\sin\theta}\frac{\partial}{\partial\varphi}\right).
\]
The operator denoted below by $\Delta_{\mathbf n}$ is the angular part of the Laplacian:
\[
\Delta_{\mathbf n}=\frac{1}{\sin\theta}\frac{\partial}{\partial\theta}
\sin\theta\frac{\partial}{\partial\theta}
+\frac{1}{\sin^2\theta}\frac{\partial^2}{\partial\varphi^2}.
\]}
\begin{equation}
\frac{\boldsymbol\nabla_{\mathbf n}}{\sqrt{j(j+1)}},
\qquad
\frac{\mathbf n\times\boldsymbol\nabla_{\mathbf n}}{\sqrt{j(j+1)}},
\qquad
\mathbf n.
\tag{7.3}
\end{equation}
We therefore define the vector spherical harmonics by
\begin{equation}
\begin{aligned}
\mathbf Y^{(\text{e})}_{jm}
&=\frac{1}{\sqrt{j(j+1)}}\boldsymbol\nabla_{\mathbf n}Y_{jm},
&P&=(-1)^j,\\
\mathbf Y^{(\text{m})}_{jm}
&=\mathbf n\times\mathbf Y^{(\text{e})}_{jm},
&P&=(-1)^{j+1},\\
\mathbf Y^{(\text{l})}_{jm}
&=\mathbf nY_{jm},
&P&=(-1)^j.
\end{aligned}
\tag{7.4}
\end{equation}

\SourcePage{32}{37}
The parity $P$ is shown alongside each vector. The three types are mutually orthogonal; $\mathbf Y^{(\text{l})}_{jm}$ is longitudinal, while $\mathbf Y^{(\text{e})}_{jm}$ and $\mathbf Y^{(\text{m})}_{jm}$ are transverse to $\mathbf n$.

The vector spherical harmonics can be expressed through scalar spherical functions. The harmonics $\mathbf Y^{(\text{m})}_{jm}$ contain spherical functions of only one order, $l=j$, whereas $\mathbf Y^{(\text{e})}_{jm}$ and $\mathbf Y^{(\text{l})}_{jm}$ contain orders $l=j\pm1$. This is evident on comparing their parities in~(7.4) with the parity $(-1)^{l+1}$ of a vector field expressed through spherical functions of order $l$.

Within each type, the vector spherical harmonics are mutually orthogonal and normalized by
\begin{equation}
\int \mathbf Y_{jm}\mathbf Y^*_{j'm'}\,do
=\delta_{jj'}\delta_{mm'}.
\tag{7.5}
\end{equation}
For $\mathbf Y^{(\text{l})}_{jm}$, this follows immediately from the normalization of $Y_{jm}$. For $\mathbf Y^{(\text{e})}_{jm}$, the normalization integral is
\[
\frac{1}{\sqrt{j(j+1)j'(j'+1)}}
\!\int \boldsymbol\nabla_{\mathbf n}Y_{jm}
\boldsymbol\nabla_{\mathbf n}Y^*_{j'm'}\,do
=-\frac{1}{\sqrt{j(j+1)j'(j'+1)}}
\!\int Y^*_{j'm'}\Delta_{\mathbf n}Y_{jm}\,do,
\]
and since $\Delta_{\mathbf n}Y_{jm}=-j(j+1)Y_{jm}$, equation~(7.5) follows. The normalization of $\mathbf Y^{(\text{m})}_{jm}$ reduces to the same integral.

The harmonics~(7.4) could also have been obtained without the direct verification of~(7.1), solely from general transformation arguments. In the preceding section, such arguments showed that a vector function $\mathbf n\varphi$ has an angular momentum $j$ equal to the order of the spherical functions entering $\varphi$. Taking $\varphi=Y_{jm}$ also fixes the projection $m$, giving $\mathbf Y^{(\text{l})}_{jm}$ immediately. The transformation argument of \S~6 is unchanged if $\mathbf n$ in $\mathbf n\varphi$ is replaced by $\boldsymbol\nabla_{\mathbf n}$ or $\mathbf n\times\boldsymbol\nabla_{\mathbf n}$, producing the other two types.

We now return to photon wave functions. For an electric-type photon ($E_j$), the vector $\mathbf A(\mathbf k)$ must have parity $(-1)^j$. Both $\mathbf Y^{(\text{e})}_{jm}$ and $\mathbf Y^{(\text{l})}_{jm}$ have this parity, but only the former satisfies transversality. For a magnetic-type photon ($M_j$), $\mathbf A(\mathbf k)$ must have
\SourcePage{33}{38}
parity $(-1)^{j+1}$, which belongs only to $\mathbf Y^{(\text{m})}_{jm}$. Thus the wave functions of a photon with specified angular momentum $j$, projection $m$, and energy $\omega$ are
\begin{equation}
\mathbf A_{\omega jm}(\mathbf k)
=\frac{4\pi^2}{\omega^{3/2}}\delta(|\mathbf k|-\omega)
\mathbf Y_{jm}(\mathbf n),
\tag{7.6}
\end{equation}
where $\mathbf Y_{jm}$ is $\mathbf Y^{(\text{e})}_{jm}$ for an electric-type photon and $\mathbf Y^{(\text{m})}_{jm}$ for a magnetic-type photon. The factor $\delta(|\mathbf k|-\omega)$ imposes the specified energy. The normalization is
\begin{equation}
\frac{1}{(2\pi)^4}\int \omega\omega'
\mathbf A^*_{\omega'j'm'}(\mathbf k)
\mathbf A_{\omega jm}(\mathbf k)\,d^3k
=\omega\delta(\omega'-\omega)\delta_{jj'}\delta_{mm'}.
\tag{7.7}
\end{equation}

For coordinate-representation wave functions, condition~(7.7) is equivalent to\footnote{This condition is of the same type as~(2.22). The factor $\delta(\omega'-\omega)$ occurs on the right because the field, a spherical wave, is considered throughout infinite space rather than in a finite volume $V=1$.}
\begin{equation}
\begin{aligned}
\frac{1}{4\pi}\int\{&\mathbf E^*_{\omega'j'm'}(\mathbf r)
\mathbf E_{\omega jm}(\mathbf r)
+\mathbf H^*_{\omega'j'm'}(\mathbf r)
\mathbf H_{\omega jm}(\mathbf r)\}\,d^3x\\
&=\omega\delta(\omega'-\omega)\delta_{jj'}\delta_{mm'}.
\end{aligned}
\tag{7.8}
\end{equation}
Indeed, expressed through the potentials, the integral on the left is
\[
\frac{1}{2\pi}\int \mathbf A^*_{\omega'j'm'}(\mathbf r)
\mathbf A_{\omega jm}(\mathbf r)\,\omega'\omega\,d^3x.
\]
Substitute
\begin{equation}
\begin{aligned}
\mathbf A_{\omega jm}(\mathbf r)
&=\int \mathbf A_{\omega jm}(\mathbf k)e^{i\mathbf k\mathbf r}
\frac{d^3k}{(2\pi)^3},\\
\mathbf A^*_{\omega'j'm'}(\mathbf r)
&=\int \mathbf A^*_{\omega'j'm'}(\mathbf k')e^{-i\mathbf k'\mathbf r}
\frac{d^3k'}{(2\pi)^3}.
\end{aligned}
\tag{7.9}
\end{equation}
Integration over $d^3x$ then gives $(2\pi)^3\delta(\mathbf k'-\mathbf k)$, which is removed by the $d^3k'$ integration, reducing the integral to~(7.7).

Up to this point we have assumed a transverse gauge, in which the scalar potential is $\Phi=0$. Other gauges for a spherical wave may, however, be more convenient in some applications.

\SourcePage{34}{39}
An admissible transformation of the potentials in momentum space is
\[
\mathbf A\to\mathbf A+\mathbf nf(\mathbf k),
\qquad
\Phi\to\Phi+f(\mathbf k),
\]
where $f(\mathbf k)$ is arbitrary. Here choose it so that the new potentials are expressed through the same spherical functions and retain a definite parity. For an electric-type photon, these requirements restrict the potentials to
\begin{equation}
\begin{aligned}
\mathbf A^{(\text{e})}_{\omega jm}(\mathbf k)
&=\frac{4\pi^2}{\omega^{3/2}}\delta(|\mathbf k|-\omega)
\left(\mathbf Y^{(\text{e})}_{jm}+C\mathbf nY_{jm}\right),\\
\Phi^{(\text{e})}_{\omega jm}(\mathbf k)
&=\frac{4\pi^2}{\omega^{3/2}}\delta(|\mathbf k|-\omega)CY_{jm},
\end{aligned}
\tag{7.10}
\end{equation}
where $C$ is arbitrary. For a magnetic-type photon, adding such a term to $\mathbf A^{(\text{m})}_{\omega jm}(\mathbf k)$ would destroy definite parity; under the same requirements, the choice~(7.6) is therefore unique.

By~(3.5) and~(7.6), the probability that a photon of definite angular momentum and parity will be detected moving in direction $\mathbf n$ within the solid-angle element $do$ is
\begin{equation}
w(\mathbf n)\,do=|\mathbf Y^{(\text{e})}_{jm}|^2\,do.
\tag{7.11}
\end{equation}
This expression was written for an $E$-type photon. Since $|\mathbf Y^{(\text{m})}_{jm}|^2=|\mathbf Y^{(\text{e})}_{jm}|^2$, however, the distributions $w(\mathbf n)$ are identical for both types.

The modulus squared $|\mathbf Y^{(\text{e})}_{jm}|^2$ is independent of the azimuth $\varphi$, because the factors $e^{\pm im\varphi}$ in the spherical functions cancel. Thus $w(\mathbf n)$ is symmetric about the $z$ axis. Moreover, each vector spherical harmonic has a definite parity, so its modulus squared is even under inversion, or under $\theta\to\pi-\theta$. Hence the expansion of $w(\theta)$ in Legendre polynomials contains only even orders. Determining its coefficients reduces to evaluating integrals of products of three spherical functions and then summing over components. Formulas from III, \S\S~107--108 perform both operations and give
\SourcePage{35}{40}
\begin{equation}
\begin{aligned}
w(\theta)={}&(-1)^{m+1}\frac{(2j+1)^2}{4\pi}
\sum_{n=0}^{\infty}(4n+1)
\begin{pmatrix}j&j&2n\\0&0&0\end{pmatrix}
\begin{pmatrix}j&j&2n\\m&-m&0\end{pmatrix}\\
&\times
\begin{Bmatrix}j&j&1\\j&j&2n\end{Bmatrix}
P_{2n}(\cos\theta).
\end{aligned}
\tag{7.12}
\end{equation}

Finally, we give the components of the vector spherical harmonics as expansions in scalar spherical functions. We use the ``spherical components'' of a vector defined in III, \S~107. For a vector $\mathbf f$, they are
\begin{equation}
f_0=if_z,
\qquad
f_{+1}=-\frac{i}{\sqrt2}(f_x+if_y),
\qquad
f_{-1}=\frac{i}{\sqrt2}(f_x-if_y).
\tag{7.13}
\end{equation}
Introduce the ``circular unit vectors''
\begin{equation}
\mathbf e^{(0)}=i\mathbf e^{(z)},
\qquad
\mathbf e^{(+1)}=-\frac{i}{\sqrt2}(\mathbf e^{(x)}+i\mathbf e^{(y)}),
\qquad
\mathbf e^{(-1)}=\frac{i}{\sqrt2}(\mathbf e^{(x)}-i\mathbf e^{(y)})
\tag{7.14}
\end{equation}
($\mathbf e^{(x,y,z)}$ are unit vectors along the $x$, $y$, and $z$ axes). Then
\begin{equation}
\mathbf f=\sum_\lambda(-1)^{1-\lambda}f_{-\lambda}\mathbf e^{(\lambda)},
\qquad
f_\lambda=(-1)^{1-\lambda}\mathbf f\mathbf e^{(-\lambda)*}
=\mathbf f\mathbf e^{(\lambda)}.
\tag{7.15}
\end{equation}
The spherical components of the vector spherical harmonics are expressed through spherical functions by means of $3j$ symbols:
\begin{equation}
\begin{aligned}
(-1)^{j+m+\lambda+1}(\mathbf Y^{(\text{e})}_{jm})_\lambda
={}&-\sqrt j
\begin{pmatrix}j+1&1&j\\m+\lambda&-\lambda&-m\end{pmatrix}
Y_{j+1,m+\lambda}\\
&+\sqrt{j+1}
\begin{pmatrix}j-1&1&j\\m+\lambda&-\lambda&-m\end{pmatrix}
Y_{j-1,m+\lambda},\\
(-1)^{j+m+\lambda+1}(\mathbf Y^{(\text{m})}_{jm})_\lambda
={}&-\sqrt{2j+1}
\begin{pmatrix}j&1&j\\m+\lambda&-\lambda&-m\end{pmatrix}
Y_{j,m+\lambda},\\
(-1)^{j+m+\lambda+1}(\mathbf Y^{(\text{l})}_{jm})_\lambda
={}&\sqrt{j+1}
\begin{pmatrix}j+1&1&j\\m+\lambda&-\lambda&-m\end{pmatrix}
Y_{j+1,m+\lambda}\\
&+\sqrt j
\begin{pmatrix}j-1&1&j\\m+\lambda&-\lambda&-m\end{pmatrix}
Y_{j-1,m+\lambda}.
\end{aligned}
\tag{7.16}
\end{equation}

These formulas are derived as follows. Each of the three vector spherical harmonics has the form $\mathbf Y_{jm}=\mathbf aY_{jm}$, where $\mathbf a$ is one of the vectors in~(7.3). Therefore
\[
\mathbf Y_{jm}=\sum_{lm'}\langle lm'|\mathbf a|jm\rangle Y_{lm'},
\]
and the problem reduces to finding matrix elements of $\mathbf a$ between eigenfunctions of orbital angular momentum. According to~(107.6) (see III),
\[
\langle lm'|a_\lambda|jm\rangle
=i(-1)^{j_{\max}-m'}
\begin{pmatrix}l&1&j\\-m'&\lambda&m\end{pmatrix}
\langle l\|\mathbf a\|j\rangle,
\]
\SourcePage{36}{41}
\noindent
where $j_{\max}$ is the greater of $l$ and $j$. It is therefore enough to know the nonzero reduced matrix elements $\langle l\|\mathbf a\|j\rangle$. They are
\begin{equation}
\begin{aligned}
\langle l-1\|\mathbf n\|l\rangle
&=\langle l\|\mathbf n\|l-1\rangle^*=i\sqrt l,\\
\langle l\|\boldsymbol\nabla_{\mathbf n}\|l-1\rangle
&=i(l-1)\sqrt l,\\
\langle l-1\|\boldsymbol\nabla_{\mathbf n}\|l\rangle
&=i(l+1)\sqrt l,\\
\langle l\|\mathbf n\times\boldsymbol\nabla_{\mathbf n}\|l\rangle
&=i\sqrt{l(l+1)(2l+1)}.
\end{aligned}
\tag{7.17}
\end{equation}
